% Compilation : pdflatex, bibtex, puis deux nouveaux passages de pdflatex.
\documentclass[11pt,a4paper]{article}

\usepackage[utf8]{inputenc}
\usepackage[T1]{fontenc}
\usepackage[french]{babel}
\usepackage{lmodern}
\usepackage{microtype}
\usepackage{geometry}
\usepackage{amsmath,amssymb}
\usepackage{siunitx}
\usepackage[version=4]{mhchem}
\usepackage{booktabs}
\usepackage{array}
\usepackage{xcolor}
\usepackage{listings}
\usepackage{enumitem}
\usepackage{hyperref}

\geometry{margin=2.5cm}
\sisetup{locale=FR,per-mode=symbol}
\DeclareSIUnit\atm{atm}
\DeclareSIUnit\molecule{molécule}
\DeclareSIUnit\ppm{ppm}
\hypersetup{
  colorlinks=true,
  linkcolor=blue!45!black,
  citecolor=blue!45!black,
  urlcolor=blue!55!black,
  pdftitle={Modèle spectral d'émissivité atmosphérique fondé sur HITRAN}
}

\definecolor{bleuclair}{RGB}{234,243,250}
\definecolor{griscode}{RGB}{247,247,247}
\newcommand{\HITRAN}{\textsc{Hitran}}
\newcommand{\HAPI}{\textsc{Hapi}}
\newcommand{\todo}[1]{%
  \par\medskip
  \noindent\fcolorbox{blue!45!black}{bleuclair}{%
    \parbox{0.94\linewidth}{\textbf{À compléter :} #1}}%
  \par\medskip
}

\lstset{
  language=Python,
  basicstyle=\ttfamily\small,
  backgroundcolor=\color{griscode},
  frame=single,
  rulecolor=\color{black!20},
  breaklines=true,
  showstringspaces=false,
  keywordstyle=\color{blue!60!black},
  commentstyle=\color{green!35!black},
  numbers=left,
  numberstyle=\tiny\color{black!55},
  xleftmargin=1.2em,
  framexleftmargin=1em
}

\title{Modèle spectral d'émissivité atmosphérique\\fondé sur la base \HITRAN}
\author{Projet climat 2026}
\date{Document de travail --- \today}

\begin{document}

\maketitle

\begin{abstract}
Ce document décrit la construction du modèle numérique qui calcule
l'absorptivité spectrale, assimilée à l'émissivité spectrale en équilibre
thermodynamique local, de couches atmosphériques contenant du
\ce{CO2}, de la vapeur d'eau, du \ce{CH4} et de l'\ce{O3}. Les sections
efficaces sont calculées raie par raie à partir de la base spectroscopique
\HITRAN, par l'intermédiaire de l'interface Python \HAPI. Elles sont combinées
aux profils verticaux de pression, de température et de fraction molaire au
moyen de la loi de Beer--Lambert. Le calcul couvre actuellement les longueurs
d'onde de \SI{4}{\micro\metre} à \SI{100}{\micro\metre} et les altitudes de
\SI{0}{\kilo\metre} à \SI{100}{\kilo\metre}.
\end{abstract}

\tableofcontents
\newpage

\section{Objectif et périmètre du modèle}

L'objectif est d'obtenir, pour une couche atmosphérique d'épaisseur
$\Delta z$, une grandeur spectrale $\varepsilon_\lambda$ dépendant de la
longueur d'onde, de l'altitude, de la pression, de la température et de la
composition gazeuse. Contrairement au premier modèle empirique, les bandes
d'absorption ne sont pas représentées par des fonctions analytiques ajustées :
elles sont reconstruites à partir des paramètres de raies de \HITRAN.

Le calcul actuel représente :

\begin{itemize}[itemsep=2pt]
  \item quatre absorbeurs : \ce{H2O}, \ce{CO2}, \ce{O3} et \ce{CH4} ;
  \item une propagation verticale, sans diffusion ;
  \item des couches homogènes et isothermes de \SI{1}{\kilo\metre} ;
  \item une grille de sortie comprise entre \SI{4}{\micro\metre} et
        \SI{100}{\micro\metre}, par pas de \SI{0.1}{\micro\metre} ;
  \item un profil atmosphérique compris entre \SI{0}{\kilo\metre} et
        \SI{100}{\kilo\metre}, par pas de \SI{1}{\kilo\metre}.
\end{itemize}

Le résultat produit par le programme est d'abord l'émissivité d'une
\emph{couche individuelle}. Le calcul du rayonnement sortant de toute la
colonne, qui doit pondérer l'émission de chaque couche par sa température et
par la transmission des couches situées au-dessus, constitue une étape
ultérieure distincte.

\section{Organisation des programmes}

Le dossier \path{modèle2_hitran_emissivité} contient quatre éléments principaux.

\begin{table}[ht]
  \centering
  \small
  \begin{tabular}{p{0.36\linewidth}p{0.56\linewidth}}
    \toprule
    Fichier & Rôle \\
    \midrule
    \path{1_section_efficace_CO2_15um.py}
      & Prototype limité à la bande du \ce{CO2} autour de
        \SI{15}{\micro\metre}. \\
    \path{2_sections_efficaces_gaz_hitran.py}
      & Téléchargement des raies, calcul Voigt des sections efficaces des
        quatre gaz et interpolation sur une grille de longueurs d'onde. \\
    \path{code_final_emissivite_avec_sections_hitran.py}
      & Calcul des densités moléculaires, des épaisseurs optiques et de
        l'émissivité d'une couche. \\
    \path{generer_table_emissivite_hitran.py}
      & Parcours des altitudes et des longueurs d'onde, puis écriture du
        tableau CSV final. \\
    \path{donnees_hitran/}
      & Cache local des tables de raies téléchargées par \HAPI. \\
    \bottomrule
  \end{tabular}
  \caption{Architecture du modèle HITRAN.}
  \label{tab:architecture}
\end{table}

La chaîne de calcul peut être résumée ainsi :

\begin{equation*}
\boxed{\text{raies \HITRAN}}
\longrightarrow
\boxed{\sigma_g(\widetilde\nu,T,p)}
\longrightarrow
\boxed{n_g(z)}
\longrightarrow
\boxed{\tau_g(\lambda,z)}
\longrightarrow
\boxed{\varepsilon_\lambda(z)}.
\end{equation*}

\subsection{Environnement et reproductibilité}

L'environnement virtuel présent dans le dossier utilise actuellement
Python~3.12.0, NumPy~2.4.6, Matplotlib~3.11.0 et le cœur de HAPI~1.3.0.0.
Ces versions doivent être figées dans un fichier de dépendances avant la
livraison définitive du modèle. Depuis la racine du dossier d'émissivité, la
table est générée avec :

\begin{lstlisting}[language=bash,numbers=none]
cd <dossier_HITRAN>
.venv/bin/python generer_table_emissivite_hitran.py
\end{lstlisting}

Dans cette commande, \texttt{<dossier\_HITRAN>} désigne le dossier
\path{modèle2_hitran_emissivité}. Le premier lancement peut nécessiter un accès
au serveur HITRAN pour remplir \path{donnees_hitran/}. Les exécutions suivantes
réutilisent les tables locales.

\section{Données spectroscopiques HITRAN}

\subsection{Base de données et molécules retenues}

\HITRAN est une base de données spectroscopiques qui fournit notamment la
position, l'intensité et les paramètres d'élargissement des raies moléculaires
\cite{gordon2022hitran}. Le programme utilise les identifiants moléculaires et
isotopologiques indiqués dans le tableau~\ref{tab:molecules}.

\begin{table}[ht]
  \centering
  \begin{tabular}{lccc}
    \toprule
    Molécule & Nom de table & Identifiant \HITRAN & Isotopologue \\
    \midrule
    \ce{H2O} & \texttt{H2O} & 1 & 1 \\
    \ce{CO2} & \texttt{CO2} & 2 & 1 \\
    \ce{O3}  & \texttt{O3}  & 3 & 1 \\
    \ce{CH4} & \texttt{CH4} & 6 & 1 \\
    \bottomrule
  \end{tabular}
  \caption{Tables demandées à \HAPI. Dans la version actuelle, seul
  l'isotopologue principal de chaque molécule est inclus.}
  \label{tab:molecules}
\end{table}

Les fonctions \texttt{db\_begin} et \texttt{fetch} de \HAPI définissent le
cache local et téléchargent les tables absentes. La présence des deux fichiers
\path{.data} et \path{.header} empêche un nouveau téléchargement. Ce contrôle
ne vérifie cependant pas que les fichiers déjà présents couvrent exactement
la plage spectrale demandée.

\subsection{Passage de la longueur d'onde au nombre d'onde}

\HAPI travaille ici en nombre d'onde $\widetilde\nu$, exprimé en
\si{\per\centi\metre}. Pour une longueur d'onde $\lambda$ exprimée en mètres,

\begin{equation}
  \widetilde\nu\;(\si{\per\centi\metre})
  = \frac{1}{100\,\lambda\;(\si{\metre})}.
  \label{eq:conversion_nombre_onde}
\end{equation}

Ainsi, l'intervalle de \SI{4}{\micro\metre} à \SI{100}{\micro\metre}
correspond à l'intervalle de \SI{2500}{\per\centi\metre} à
\SI{100}{\per\centi\metre}. Les calculs \HAPI sont effectués sur une grille
uniforme en nombre d'onde, avec un pas de
\SI{0.05}{\per\centi\metre}.

\subsection{Profil de raie de Voigt}

La fonction \texttt{absorptionCoefficient\_Voigt} de \HAPI calcule une section
efficace spectrale en combinant l'élargissement Doppler, de forme gaussienne,
et l'élargissement collisionnel, de forme lorentzienne
\cite{kochanov2016hapi}. De manière symbolique,

\begin{equation}
  \sigma_g(\widetilde\nu,T,p)
  = \sum_{j\in g} S_j(T)\,
    V\!\left(\widetilde\nu-\widetilde\nu_j;
      \gamma_{D,j}(T),\gamma_{L,j}(T,p)\right),
  \label{eq:section_voigt}
\end{equation}

où $S_j$ est l'intensité de la raie $j$, $\widetilde\nu_j$ sa position et
$V$ le profil de Voigt normalisé. La température modifie notamment les
intensités et l'élargissement Doppler ; la pression intervient dans
l'élargissement collisionnel.

L'option \texttt{HITRAN\_units=True} renvoie $\sigma_g$ en
\si{\centi\metre\squared\per\molecule}. Le programme applique donc

\begin{equation}
  \sigma_g\;(\si{\metre\squared\per\molecule})
  = 10^{-4}\,
    \sigma_g\;(\si{\centi\metre\squared\per\molecule}).
  \label{eq:conversion_section}
\end{equation}

Pour chaque couche, la pression calculée en pascals est convertie en
atmosphères avant l'appel à \HAPI :

\begin{equation}
  p_{\mathrm{atm}} = \frac{p_{\mathrm{Pa}}}{101325}.
\end{equation}

\subsection{Interpolation sur la grille de sortie}

Le spectre calculé par \HAPI est tabulé sur une grille uniforme en nombre
d'onde. La fonction \texttt{numpy.interp} évalue ensuite la section efficace
aux nombres d'onde associés à la grille uniforme en longueur d'onde. Hors de
la plage $[100,2500]$~\si{\per\centi\metre}, la section est fixée à zéro.

Il faut distinguer les deux résolutions : le calcul interne a un pas de
\SI{0.05}{\per\centi\metre}, tandis que le CSV est échantillonné tous les
\SI{0.1}{\micro\metre}. Un pas uniforme en longueur d'onde n'est pas uniforme
en nombre d'onde ; la résolution spectrale effective du tableau varie donc
avec $\lambda$.

\section{Profil vertical de l'atmosphère}

Les propriétés atmosphériques sont fournies par la fonction
\texttt{get\_gases\_ppm} du fichier
\path{températures_et_concentrations/code_H20_CH4_CO2_O3.py}.

\subsection{Température et pression}

La température suit un profil affine par morceaux, dont les limites sont
0, 11, 20, 32, 50, 71 et \SI{85}{\kilo\metre}. Dans chaque couche, on écrit

\begin{equation}
  T(z)=T_0+\alpha(z-z_0),
\end{equation}

où $\alpha$ est le gradient thermique local. La pression est obtenue à partir
de l'équilibre hydrostatique et de la loi des gaz parfaits. Pour
$\alpha\neq0$,

\begin{equation}
  p(z)=p_0\left(\frac{T(z)}{T_0}\right)^{-gM/(R\alpha)},
  \label{eq:pression_gradient}
\end{equation}

et, pour une couche isotherme,

\begin{equation}
  p(z)=p_0\exp\!\left[-\frac{gM(z-z_0)}{RT_0}\right].
  \label{eq:pression_isotherme}
\end{equation}

Le modèle prend $g=\SI{9.80665}{\metre\per\second\squared}$,
$M=\SI{0.0289644}{\kilogram\per\mole}$ et
$R=\SI{8.31445}{\joule\per\mole\per\kelvin}$.

\subsection{Fractions molaires des gaz}

Les concentrations renvoyées sous le nom de « proportion » sont exprimées en
parties par million. Leur conversion en fraction molaire est

\begin{equation}
  x_g = C_g\;(\mathrm{ppm})\times10^{-6}.
  \label{eq:ppm}
\end{equation}

Les profils actuellement implémentés sont les suivants :

\begin{itemize}
  \item \ce{CO2} : fraction molaire constante de \SI{425}{\ppm} ;
  \item \ce{CH4} : \SI{1.9}{\ppm} jusqu'à \SI{11}{\kilo\metre}, puis
    décroissance exponentielle de hauteur caractéristique
    \SI{25}{\kilo\metre} ;
  \item \ce{O3} : fond de \SI{0.1}{\ppm}, remplacé entre 15 et
    \SI{35}{\kilo\metre} par un profil gaussien centré à
    \SI{25}{\kilo\metre}, de maximum \SI{8}{\ppm} ;
  \item \ce{H2O} : calcul à partir d'une humidité relative décroissante dans
    la troposphère et de la pression de vapeur saturante, avec un plancher de
    \SI{4}{\ppm} au-dessus de la tropopause.
\end{itemize}

Ces profils sont des paramétrisations simplifiées et non des observations
météorologiques instantanées. Leur rôle est de fournir une atmosphère de
référence cohérente pour tester la chaîne radiative.

\section{De la section efficace à l'émissivité}

\subsection{Densité moléculaire}

La densité moléculaire totale de l'air est déduite de la loi des gaz parfaits
sous forme microscopique :

\begin{equation}
  n_{\mathrm{air}}(z)=\frac{p(z)}{k_{\mathrm B}T(z)},
  \qquad
  k_{\mathrm B}=\SI{1.380649e-23}{\joule\per\kelvin}.
  \label{eq:densite_air}
\end{equation}

La densité du gaz $g$ vaut alors

\begin{equation}
  n_g(z)=x_g(z)n_{\mathrm{air}}(z).
  \label{eq:densite_gaz}
\end{equation}

Les unités sont cohérentes : $p/(k_{\mathrm B}T)$ est en
\si{\per\metre\cubed} et $x_g$ est sans dimension.

\subsection{Épaisseur optique de chaque gaz}

Pour une couche homogène d'épaisseur géométrique $\Delta z$, l'épaisseur
optique monochromatique due au gaz $g$ est

\begin{equation}
  \tau_g(\lambda,z)
  = \sigma_g(\lambda,T(z),p(z))\,n_g(z)\,\Delta z.
  \label{eq:tau_gaz}
\end{equation}

L'épaisseur optique est sans dimension. En supposant que les absorptions des
différents gaz sont indépendantes, elles s'additionnent :

\begin{equation}
  \tau_{\mathrm{tot}}
  = \tau_{\ce{CO2}}+\tau_{\ce{H2O}}
  +\tau_{\ce{CH4}}+\tau_{\ce{O3}}.
  \label{eq:tau_total}
\end{equation}

\subsection{Transmission, absorptivité et émissivité}

La loi de Beer--Lambert donne la transmission spectrale de la couche :

\begin{equation}
  \mathcal{T}_\lambda=\exp(-\tau_{\mathrm{tot}}).
  \label{eq:transmission}
\end{equation}

En l'absence de diffusion et de réflexion, son absorptivité est

\begin{equation}
  \alpha_\lambda=1-\mathcal{T}_\lambda
  =1-\exp(-\tau_{\mathrm{tot}}).
  \label{eq:absorptivite}
\end{equation}

Sous l'hypothèse d'équilibre thermodynamique local, la loi de Kirchhoff permet
d'identifier l'émissivité spectrale à l'absorptivité spectrale :

\begin{equation}
  \boxed{\varepsilon_\lambda
  =1-\exp(-\tau_{\mathrm{tot}})}.
  \label{eq:emissivite}
\end{equation}

Cette quantité décrit la capacité d'émission d'une couche à une longueur
d'onde donnée. La luminance effectivement émise dans une direction ferait en
plus intervenir la loi de Planck $B_\lambda(T)$ et la géométrie du trajet.

\section{Génération du tableau numérique}

Pour chaque altitude entière $z$ entre 0 et \SI{100}{\kilo\metre}, le
programme suit les étapes ci-dessous :

\begin{enumerate}
  \item calcul de $T(z)$, $p(z)$ et des concentrations des quatre gaz ;
  \item calcul par \HAPI des quatre spectres $\sigma_g$ aux valeurs locales de
        température et de pression ;
  \item interpolation de ces spectres sur 961 longueurs d'onde, de 4 à
        \SI{100}{\micro\metre} ;
  \item calcul de $n_g$, $\tau_g$ et $\varepsilon_\lambda$ pour une couche de
        \SI{1000}{\metre} ;
  \item écriture d'une ligne par couple $(z,\lambda)$ dans le fichier CSV.
\end{enumerate}

La grille complète contient $101\times961=97061$ lignes de données, auxquelles
s'ajoute la ligne d'en-tête. Les colonnes enregistrent :

\begin{itemize}[itemsep=2pt]
  \item l'altitude, la longueur d'onde et le nombre d'onde ;
  \item la pression, l'épaisseur de couche et la température ;
  \item les quatre concentrations en ppm ;
  \item les quatre sections efficaces en
        \si{\metre\squared\per\molecule} ;
  \item les quatre épaisseurs optiques ;
  \item l'émissivité spectrale totale de la couche.
\end{itemize}

L'épaisseur optique totale n'est pas écrite directement, mais peut être
reconstruite en sommant les quatre colonnes $\tau_g$ ou à partir de
$\tau_{\mathrm{tot}}=-\ln(1-\varepsilon)$ lorsque $\varepsilon<1$.

\section{Vérifications et tests de cohérence}

Une validation minimale du calcul doit comprendre les tests suivants.

\subsection{Contrôles d'unités et de bornes}

\begin{itemize}
  \item vérifier $\sigma_g\geq0$, $\tau_g\geq0$ et
        $0\leq\varepsilon_\lambda\leq1$ ;
  \item vérifier la correspondance
        \SI{4}{\micro\metre}$\leftrightarrow$
        \SI{2500}{\per\centi\metre} et
        \SI{100}{\micro\metre}$\leftrightarrow$
        \SI{100}{\per\centi\metre} ;
  \item vérifier qu'une concentration nulle ou une section nulle donne
        $\tau_g=0$ ;
  \item vérifier que $\varepsilon_\lambda\simeq\tau_{\mathrm{tot}}$ dans la
        limite optiquement mince $\tau_{\mathrm{tot}}\ll1$ ;
  \item vérifier que $\varepsilon_\lambda\to1$ lorsque
        $\tau_{\mathrm{tot}}\gg1$.
\end{itemize}

\subsection{Contrôle du fichier produit}

Une génération complète doit contenir exactement 97061 lignes de données,
101 altitudes distinctes et 961 longueurs d'onde par altitude. L'écriture est
effectuée altitude après altitude et le fichier est vidé sur disque après
chaque bloc. Une interruption laisse donc un CSV lisible mais incomplet ; le
nombre de lignes et la dernière altitude doivent être contrôlés avant toute
exploitation.

\subsection{Comparaison au modèle empirique}

Le script
\path{comparaison_modeles_sections_efficaces/comparer_sections_efficaces.py}
compare, gaz par gaz, les
sections efficaces empiriques et celles de \HITRAN à
$T=\SI{296}{\kelvin}$ et $p=\SI{1}{\atm}$. Cette comparaison doit porter sur :

\begin{itemize}
  \item la position des principales bandes d'absorption ;
  \item l'ordre de grandeur des sections efficaces ;
  \item la structure fine des raies, présente dans \HITRAN mais absente des
        fonctions empiriques ;
  \item l'effet du sous-échantillonnage lors du passage au pas de
        \SI{0.1}{\micro\metre}.
\end{itemize}

\todo{Insérer les quatre figures de comparaison et commenter quantitativement
les écarts dans les bandes principales de \ce{CO2}, \ce{H2O}, \ce{CH4} et
\ce{O3}.}

\section{Hypothèses, limites et précautions d'interprétation}

Les résultats doivent être interprétés à la lumière des hypothèses suivantes.

\begin{enumerate}
  \item \textbf{Isotopologues.} Seul l'isotopologue d'identifiant 1 est
  téléchargé pour chaque molécule. L'absorption des isotopologues minoritaires
  est donc absente.

  \item \textbf{Profil de Voigt.} Le profil Voigt est une approximation
  standard, mais les effets de mélange de raies, de continuum, de dépendance
  en vitesse et les choix de coupure des ailes ne sont pas étudiés ici.

  \item \textbf{Milieu homogène.} Chaque tranche de \SI{1}{\kilo\metre} est
  représentée par une seule pression, une seule température et une seule
  composition. Le point exact auquel ces propriétés représentent la couche
  (base, centre ou moyenne) doit être fixé explicitement dans la suite du
  projet.

  \item \textbf{Absence de diffusion.} La diffusion moléculaire, les aérosols,
  les nuages et les réflexions sont négligés. L'assimilation
  $\varepsilon_\lambda=1-\mathcal{T}_\lambda$ repose sur cette hypothèse.

  \item \textbf{Équilibre thermodynamique local.} L'égalité entre
  absorptivité et émissivité est utilisée couche par couche.

  \item \textbf{Géométrie verticale.} L'équation~\eqref{eq:tau_gaz} utilise
  une longueur de trajet verticale $\Delta z$. Un rayon oblique traverse une
  distance plus grande et requiert un facteur géométrique adapté.

  \item \textbf{Émissivité d'une colonne non isotherme.} On peut sommer les
  épaisseurs optiques pour obtenir la transmission totale. En revanche, on ne
  peut pas attribuer une température unique à toute la colonne ni sommer
  directement les émissions des couches : chacune doit être pondérée par la
  loi de Planck et atténuée par les couches sus-jacentes.

  \item \textbf{Profils atmosphériques simplifiés.} Les profils de gaz,
  notamment ceux de la vapeur d'eau et de l'ozone, sont paramétriques. Le
  profil thermique au-dessus de \SI{85}{\kilo\metre} est également une
  extrapolation simplifiée.

  \item \textbf{Échantillonnage.} Le pas de sortie de
  \SI{0.1}{\micro\metre} peut masquer des raies beaucoup plus étroites. Une
  moyenne sur une bande instrumentale n'est pas équivalente à une simple
  interpolation ponctuelle.
\end{enumerate}

\section{Pistes d'amélioration}

Les développements prioritaires sont :

\begin{itemize}
  \item définir les propriétés de chaque couche à son altitude médiane ;
  \item inclure les isotopologues pertinents et documenter leur abondance ;
  \item comparer plusieurs résolutions spectrales et effectuer, si nécessaire,
        une moyenne de bande conservant la transmission ;
  \item automatiser les tests de complétude et les tests physiques du CSV ;
  \item construire l'équation de transfert radiatif multicouche avec la loi de
        Planck ;
  \item comparer le rayonnement sortant à un spectre de référence ou à une
        simulation reconnue ;
  \item étudier la sensibilité aux profils de température, de pression et de
        concentration.
\end{itemize}

\todo{Ajouter une section de validation quantitative : cas tests, graphiques,
temps de calcul, convergence en résolution spectrale et bilan d'erreur.}

\section{Conclusion provisoire}

Le modèle relie de façon explicite les données spectroscopiques de \HITRAN à
une propriété radiative de couche. Sa chaîne physique est cohérente : les
sections efficaces dépendant de $T$ et $p$ sont multipliées par les densités
moléculaires et la longueur de trajet pour former les épaisseurs optiques, puis
la loi de Beer--Lambert donne l'absorptivité spectrale. La prochaine étape ne
consiste pas seulement à raffiner les spectres : elle doit surtout transformer
ces propriétés de couche en un véritable calcul de transfert radiatif pour une
atmosphère non isotherme.

\bibliographystyle{plain}
\bibliography{references_hitran}

\end{document}
